\section*{Chapitre \ref{ch:g10:soil-formation} --- Le sol, une ressource lente à se former}

\begin{solution}{exo:g10:soil-formation:1}
L'\omterm{def:g10:soil-formation:chemical-weathering}{altération chimique} est la transformation des minéraux d'une roche par des
réactions avec l'eau et les gaz qu'elle transporte (dioxyde de carbone,
dioxygène)~: certains minéraux se dissolvent, d'autres deviennent de
nouveaux minéraux, stables en surface. L'\omterm{def:g10:soil-formation:hydrolysis}{hydrolyse} est l'\omterm{def:g10:soil-formation:chemical-weathering}{altération
chimique} d'un silicate par l'eau et ses ions hydrogène~: les ions
métalliques passent en solution et le reste du minéral est reconstruit en un
\omterm{def:g10:soil-formation:clay}{minéral argileux}.
\end{solution}

\begin{solution}{exo:g10:soil-formation:2}
La \omterm{def:g10:soil-formation:weathering-profile}{saprolite} est une roche décomposée sur place~: tendre et terreuse, elle
garde les grains, les bandes et la structure de sa \omterm{def:g10:soil-formation:weathering-profile}{roche mère} et passe vers
le bas, sans rupture, à la roche saine. Un sol formé dans du sable de
rivière repose sur un matériau venu d'ailleurs~: ses grains sont arrondis
et triés, et il ne montre rien de la structure de la roche sous-jacente.
\end{solution}

\begin{solution}{exo:g10:soil-formation:3}
Les feldspaths deviennent de l'argile (kaolinite), en libérant dans l'eau du
potassium, du sodium ou du calcium et de la silice~; le mica noir devient de
l'argile et des oxydes de fer~; le quartz reste intact, sous forme de
sable.
\end{solution}

\begin{solution}{exo:g10:soil-formation:4}
Le fer des minéraux sombres, \ce{Fe^2+}, est oxydé par le dioxygène apporté
par l'eau en oxydes de fer(III) insolubles~: la goethite jaune-brun et
l'hématite rouge, qui enrobent les grains et colorent tout le matériau.
\end{solution}

\begin{solution}{exo:g10:soil-formation:5}
\qty{0.017}{mm} par an et \qty{1.537}{mm} par an~: $1.537/0.017 \approx 90$
fois plus.
\end{solution}

\begin{solution}{exo:g10:soil-formation:6}
Calcium~: 1 et 1. Carbone~: $1 + 1 = 2$ et 2. Oxygène~: $3 + 2 + 1 = 6$ et
$2 \times 3 = 6$. Hydrogène~: 2 et 2. Charge~: 0, et $+2 - 2 = 0$.
\end{solution}

\begin{solution}{exo:g10:soil-formation:7}
La dissolution de la calcite ne produit que des ions dissous, \ce{Ca^2+} et
\ce{HCO3-}~: tout est emporté et un calcaire pur ne laisse aucun résidu
solide. L'\omterm{def:g10:soil-formation:hydrolysis}{hydrolyse} du feldspath produit un solide, la kaolinite, qui
reste~; et le quartz du granite ne réagit presque pas. Le granite laisse donc
de l'argile et du sable.
\end{solution}

\begin{solution}{exo:g10:soil-formation:8}
Les paillettes d'argile et l'humus portent des charges négatives qui
retiennent à leur surface les ions positifs potassium \ce{K+}, si bien que
l'eau de pluie qui traverse le sol ne les emporte pas. Une racine libère des
ions hydrogène, qui prennent la place du potassium sur les paillettes~; le
potassium libéré dans l'eau du sol est absorbé par la racine.
\end{solution}

\begin{solution}{exo:g10:soil-formation:9}
L'\omterm{def:g10:soil-formation:chemical-weathering}{altération chimique} a besoin d'eau, abondante sous les climats humides et
rare sous les climats secs (et gelée sous les climats froids)~; et ses
réactions sont plus rapides quand il fait chaud. Ce n'est que sous la
chaleur et avec une eau abondante, pendant longtemps, que les réactions vont
assez loin pour retirer presque tout ce qui est soluble et laisser des
oxydes de fer et d'aluminium.
\end{solution}

\begin{solution}{exo:g10:soil-formation:10}
$10/0.017 \approx 590$ ans pour \qty{1}{cm}~; $300/0.017 \approx 17\,600$
ans pour \qty{30}{cm}~; $1000/0.017 \approx 59\,000$ ans pour \qty{1}{m}.
\end{solution}

\begin{solution}{exo:g10:soil-formation:11}
Volume~: $10\,000 \times 0.0015 = \qty{15}{m^3}$~; masse~: $15 \times 1200 =
\qty{18000}{kg}$, soit \num{18}~tonnes par hectare et par an.
\end{solution}

\begin{solution}{exo:g10:soil-formation:12}
$46\,000/0.017 \approx 2.7$~millions d'années. Une \omterm{def:g10:soil-formation:weathering-profile}{saprolite} aussi épaisse
exige un paysage resté stable, sans que son manteau d'altération soit
décapé, pendant des millions d'années.
\end{solution}

\begin{solution}{exo:g10:soil-formation:13}
Goethite~: \ce{4 Fe^2+ + O2 + 6 H2O -> 4 FeO(OH) + 8 H+}~: fer 4 et 4,
oxygène $2 + 6 = 8$ et $4 \times 2 = 8$, hydrogène 12 et $4 + 8 = 12$, charge
$+8$ et $+8$. Hématite~: \ce{4 Fe^2+ + O2 + 4 H2O -> 2 Fe2O3 + 8 H+}~: fer 4
et 4, oxygène $2 + 4 = 6$ et 6, hydrogène 8 et 8, charge $+8$ et $+8$.
\end{solution}

\begin{solution}{exo:g10:soil-formation:14}
Perte nette $0.082 - 0.017 = \qty{0.065}{mm}$ par an~; $300/0.065 \approx 4600$
ans. Sous un labour conventionnel, $300/1.52 \approx 200$ ans~: plus de
vingt fois moins.
\end{solution}

\begin{solution}{exo:g10:soil-formation:15}
La tonne contient \qty{300}{kg} d'orthose. Chaque kilogramme donne
\qty{464}{g} de kaolinite~: $300 \times 0.464 \approx \qty{139}{kg}$. Potassium~:
$300 \times 0.1405 \approx \qty{42}{kg}$.
\end{solution}

\begin{solution}{pb:g10:soil-formation:1}
\textbf{1.} $4 \times 0.3048 \approx \qty{1.22}{m}$.

\textbf{2.} $0.5 \times 30.48 \approx \qty{15}{cm}$.

\textbf{3.} L'humus, la matière organique sombre fabriquée par les
décomposeurs à partir des plantes mortes.

\textbf{4.} Des oxydes de fer, la goethite et l'hématite, formés par
l'oxydation du fer(II) des minéraux du matériau parental.

\textbf{5.} Non. En dessous, le profil passerait à une roche altérée, la
\omterm{def:g10:soil-formation:weathering-profile}{saprolite} (ou un autre matériau altéré), puis à la \omterm{def:g10:soil-formation:weathering-profile}{roche mère} saine.

\textbf{6.} \ce{2 KAlSi3O8 + 2 H+ + 9 H2O -> Al2Si2O5(OH)4 + 2 K+ + 4 H4SiO4}.

\textbf{7.} $1000/278.33 \approx \qty{3.59}{mol}$.

\textbf{8.} $3.59/2 \approx \qty{1.80}{mol}$ de kaolinite, $1.80 \times 258.16
\approx \qty{464}{g}$.

\textbf{9.} Les ions potassium et la silice dissoute~: l'eau les emporte aux
rivières, et les rivières à la mer.

\textbf{10.} Il ne réagit presque pas~: les grains de quartz restent sous
forme de sable.

\textbf{11.} $1219/0.017 \approx 72\,000$ ans.

\textbf{12.} $1.537 - 0.017 = \qty{1.52}{mm}$ par an.

\textbf{13.} $1219/1.52 \approx 800$ ans.

\textbf{14.} $72\,000/800 = 90$~: le profil se perd quatre-vingt-dix fois
plus vite qu'il ne s'est formé.

\textbf{15.} Non~: $0.013 < 0.017$, l'érosion est un peu plus lente que la
production et le sol se maintient.

\textbf{16.} $T = S/(E - P)$.

\textbf{17.} $300/1.52 \approx 200$ ans.

\textbf{18.} $300/(0.082 - 0.017) = 300/0.065 \approx 4600$ ans.

\textbf{19.} Un sol perdu en quelques siècles met des dizaines de milliers
d'années à se reformer~: à l'échelle des vies humaines et des sociétés, ce
qui est perdu n'est pas remplacé.

\textbf{20.} $1000/1.52 \approx 660$ ans~: \textbf{un mètre de sol se perd
en environ 660 ans, quelques siècles, sous un labour conventionnel}.
\end{solution}
